\section{Related Work}
\label{sec:related}

\subsection{Curvilinear Structure Segmentation}
\label{sec:related_css}

Curvilinear structure segmentation aims to delineate thin, elongated, and
branching targets such as vessels, nerves, cracks, and roads. U-Net
\cite{ronneberger2015unet} established the encoder--decoder template, and
attention-based and multi-scale variants such as CS-Net \cite{mou2019csnet}
and CS2-Net \cite{mou2021cs2net} improved the delineation of vessels and
corneal nerves across imaging modalities. Domain-specific architectures push
in-domain accuracy further: DeepCrack \cite{liu2019deepcrack} aggregates
hierarchical convolutional features for crack segmentation, HACNet~V2
\cite{chen2025hacnetv2} preserves full-resolution detail for pixel-level
crack detection, and nnWNet \cite{zhou2025nnwnet} revisits
transformer--convolution hybrids. These models obtain strong results when
training and test data share a domain, but they rely on large annotated
datasets and degrade under modality or acquisition shift.

Pixel-wise overlap does not capture connectivity, so a parallel line of work
adds topological supervision. clDice \cite{shit2021cldice} measures overlap
between masks and their skeletons and provides a soft, topology-preserving
objective; Skeleton Recall \cite{kirchhoff2024skeletonrecall} supervises the
skeleton directly to conserve connectivity at lower cost. TopoMortar
\cite{valverde2025topomortar} shows that the benefit of a topology loss
depends on dataset challenges such as small training sets, noisy labels, and
out-of-distribution test images, which motivates evaluating segmentation
methods across many datasets rather than on a single benchmark.

Annotation-efficient and universal curvilinear models reduce the supervision
requirement. UCS \cite{zhu2023ucs} adapts SAM with a sparse adapter and
FFT-based prompt generation for cross-domain curvilinear segmentation. SACM
\cite{zhu2026sacm} introduces dual-level adapters, multi-layer adapter
fusion, and a dual-stage prompt-free decoder, and trains with only 18 images.
Our work shares SACM's prompt-free, low-data setting but studies a different
question: how the decoder should deploy its candidate masks. Topology-aware
losses are training signals in our method rather than the deployment
mechanism.

\subsection{SAM and Parameter-Efficient Adaptation}
\label{sec:related_sam}

SAM \cite{kirillov2023sam} is a promptable segmentation foundation model with
a frozen ViT image encoder, a prompt encoder, and a lightweight mask decoder
that returns multiple masks together with predicted IoU scores. Full
fine-tuning is expensive and prone to overfitting on small target sets, so
parameter-efficient adaptation has become the default. SAM-Adapter
\cite{chen2023samadapter} injects task-specific adapters into the encoder;
Conv-LoRA \cite{zhong2024convlora} adds convolutional parameters to low-rank
adaptation to restore local inductive bias; CAT-SAM \cite{xiao2024catsam}
performs conditional few-shot tuning of the encoder and decoder; MedSAM
\cite{ma2024medsam} and SAM-Med2D \cite{cheng2023sammed2d} scale medical
fine-tuning; and HQ-SAM \cite{ke2023hqsam} introduces a high-quality output
token for intricate structures.

Prompt-free adaptation is particularly relevant to curvilinear structures
because point and box prompts are ill-posed for dense, branching targets.
SAM-OCTA \cite{wang2024sam_octa} applies low-rank adaptation with prompt-point
generation to OCTA vessel segmentation; CWSAM \cite{pu2025cwsam} uses
class-wise adapters for SAR imagery; and SACM \cite{zhu2026sacm} replaces
prompts with dual-level adapters, adapter fusion, and a dual-stage decoder.
DCD follows this prompt-free adapter line and keeps the frozen-SAM scaffold;
its contribution is the decoding and deployment rule rather than a new
adapter architecture.

\subsection{Multi-Candidate Decoding and Test-Time Selection}
\label{sec:related_selection}

Segmentation models often represent ambiguity with multiple outputs. SAM
returns several masks and predicted IoU scores \cite{kirillov2023sam}, and
HQ-SAM adds a quality token \cite{ke2023hqsam}. Probabilistic U-Net
\cite{kohl2018probunet} learns a distribution over masks and generates
multiple plausible hypotheses. In prompt-free curvilinear decoding, SACM
\cite{zhu2026sacm} ranks multiple candidate descriptors by Stage-1
confidence, refines them in Stage-2, and selects the final mask with an
MLP-based IoU predictor.

These approaches rely on a test-time score, selector, or fusion. The score is
learned on the Seen distribution, so its calibration can change under domain
shift; selecting for overlap may favor thicker, better-connected masks while
sacrificing thin branches; and a selector or fusion branch adds inference
cost. DCD keeps multiple candidates as training-time experts but commits to a
fixed ranked slot at inference. Gating only determines the candidate order,
and the fusion branch is discarded after training; no input-dependent
selection or fusion is deployed. This differs from mixture-of-experts gating,
where the deployed combination depends on the input, and from ensembles or
test-time augmentation, which aggregate several predictions.
